% Part: first-order-logic
% Chapter: proof-systems
% Section: seqeunt-calculus

\documentclass[../../../include/open-logic-section]{subfiles}

\begin{document}

\iftag{FOL}
      {\olfileid[pt]{fol}{prf}{seq}}
      {\olfileid[pt]{pl}{prf}{seq}}

\olsection{O Cálculo de Sequentes}

Enquanto muitos sistemas de !!{derivation} operam com disposições de
!!{sentence}s, o cálculo de sequentes opera com \emph{sequentes}. Um
sequente é uma expressão da forma
\[
!A_1, \dots, !A_m \Sequent !B_1, \dots, !B_m,
\]
isto é, um par de sequências de !!{sentence}s, separadas pelo símbolo de
sequente~$\Sequent$. Qualquer uma das sequências pode ser vazia.  Uma
!!{derivation} no cálculo de sequentes é uma árvore de sequentes, em que
os sequentes mais altos são de uma forma especial (são chamados de
``sequentes iniciais'' ou ``axiomas'') e todo outro sequente segue dos
sequentes imediatamente acima dele por uma das regras de inferência. As
regras de inferência ou manipulam as !!{sentence}s nos sequentes
(adicionando, removendo ou rearranjando-as à esquerda ou à direita), ou
introduzem uma !!{formula} complexa na conclusão da regra.
Por exemplo, a regra $\LeftR{\land}$ permite a inferência de $!A,
\Gamma \Sequent \Delta$ para $!A \land !B, \Gamma \Sequent \Delta$, e
a $\RightR{\lif}$ permite a inferência de $!A, \Gamma \Sequent
\Delta, !B$ para $\Gamma \Sequent \Delta, !A \lif !B$, para quaisquer
$\Gamma$, $\Delta$, $!A$ e~$!B$. (Em particular, $\Gamma$ e~$\Delta$
podem ser vazios.)

A relação $\Proves$ baseada no cálculo de sequentes é definida da
seguinte forma: $\Gamma \Proves !A$ se, e somente se, há alguma
sequência $\Gamma_0$ tal que toda $!A$ em $\Gamma_0$ está em~$\Gamma$ e
há uma !!{derivation} com o sequente~$\Gamma_0 \Sequent !A$ em sua raiz.
$!A$ é um teorema no cálculo de sequentes se o sequente~$\Sequent !A$
tem uma !!{derivation}. Por exemplo, aqui está uma !!{derivation} que
mostra que $\Proves (!A \land !B) \lif !A$:
\begin{prooftree}
  \Axiom$!A \fCenter !A$
  \RightLabel{\LeftR{\land}}
  \UnaryInf$!A \land !B \fCenter !A$
  \RightLabel{\RightR{\lif}}
  \UnaryInf$\fCenter (!A \land !B) \lif !A$
\end{prooftree}

Um conjunto $\Gamma$ é inconsistente no cálculo de sequentes se há
uma !!{derivation} de $\Gamma_0 \Sequent$ (em que toda $!A \in \Gamma_0$
está em~$\Gamma$ e o lado direito do sequente é vazio).  Usando a
regra \RightR{\Weakening}, qualquer !!{sentence} pode ser !!{derive}da
a partir de um conjunto inconsistente.

O cálculo de sequentes foi inventado na década de 1930 por Gerhard
Gentzen. Por causa de seu desenho sistemático e simétrico, é um
formalismo muito útil para desenvolver uma teoria de !!{derivation}s. É
relativamente fácil encontrar !!{derivation}s no cálculo de sequentes,
mas essas !!{derivation}s são frequentemente difíceis de ler e sua
conexão com provas às vezes não é fácil de ver. Provou ser uma abordagem
muito elegante a sistemas de !!{derivation}, no entanto, e muitas
lógicas têm sistemas de cálculo de sequentes.

\end{document}
